% Fragmento público generado desde propietarios íntegros.
% Residencia: 52.6.2.
% Fuente: ../incoming/radion_monografia_20260827/manuscrito/sections/07_incertidumbre_oscilador.tex:156-197
\paragraph{Échange énergétique entre les secteurs électrique et magnétique}

Les deux termes de \eqref{eq:HLC} deviennent
\begin{equation}
U_E(\theta)=\hret\omega\cos^2\theta,
\qquad
U_B(\theta)=\hret\omega\sin^2\theta.
\label{eq:energy-exchange}
\end{equation}
Par conséquent,
\[
U_E+U_B=\hret\omega.
\]
La polarité électrique--magnétique ne consiste pas à étiqueter arbitrairement
les foyers. Elle consiste en deux réserves conjuguées qui s'échangent pendant
l'orbite tout en maintenant constantes la fréquence et l'action totale.

Le quotient commun se prolonge :
\begin{equation}
\boxed{
e^{-\eta}
=\frac{b_S}{a_S}
=\frac{\Delta P}{\Delta Q}
=\frac{Z_\eta}{Z_*}
=\sqrt{\frac{A-C^*}{A+C^*}}.}
\label{eq:shape-chain-LC}
\end{equation}
Les produits conservés sont
\[
a_Sb_S=2\hret,\qquad
\det V_\eta=\frac{\hret^2}{4},
\qquad
L_\eta C_\eta=\omega^{-2}.
\]

\begin{exactbox}
L'oscillateur LC est la réalisation physique minimale de la double lecture : le
produit conserve l'action et la fréquence ; le quotient conserve l'orientation et
l'impédance. Un cycle complet échange électricité et magnétisme sans
perdre l'aire \(\oint P\,\diff Q=\hfull\).
\end{exactbox}

